% ═════════════════════════════════════════════════════════════════════════════
\section{Deductive and Proof-Based Methods}
% ═════════════════════════════════════════════════════════════════════════════

Deductive approaches shift the verification task from automated state-space
exploration to the construction of formal proofs, often assisted by theorem
provers or proof assistants. The idea is to find invariants, properties
that hold at every reachable state of the system for every~$n$, and prove
that they are preserved by every system transition, regardless of the number
of processes~\cite{OG76,Lam94}. Since the proof is carried out symbolically
over an arbitrary~$n$, this style of reasoning handles unbounded parameters
without requiring enumeration of instances.

Classic work by Owicki and Gries~\cite{OG76} laid the groundwork for
assertional reasoning about concurrent programs: each process is annotated
with pre- and post-conditions, and an interference freedom check ensures that
no process's execution invalidates another's assertions. The method was
designed for a fixed, known number of processes, but the underlying proof
discipline, inductive invariants plus non-interference, scales naturally
to the parameterized case by carrying out the annotation and check for a
symbolic, arbitrary process in a family. Lamport's TLA+~\cite{Lam94} provides
a formal language for specifying and proving properties of concurrent and
distributed systems in this spirit, and has been applied to industrial-scale
protocols at Amazon Web Services and Intel, among others.

A significant challenge for deductive methods is the discovery of inductive
invariants: properties strong enough to be preserved under all transitions and
to imply the desired safety property, but not so strong as to be unprovable.
The \emph{invisible invariants} technique of Pnueli, Ruah, and
Zuck~\cite{PRZ01} addresses this by automatically generating candidate
invariants from small finite instances of the system and then checking whether
they hold inductively for all~$n$. The idea is that the right invariant is
often already visible in small instances and can be generalized to the full
parameterized family. This makes deductive verification partly automatic, at
the cost of completeness: if the right invariant is not among those generated,
the procedure does not terminate with a proof.

More recent work integrates deductive reasoning with interactive proof
assistants, allowing arbitrarily complex invariants to be stated and verified
with machine-checked proofs. The Ivy system~\cite{PLSS16}, which we return to
in Section~7 in its SMT-based incarnation, is representative of this direction:
it combines an interactive specification language with automated checking of
proof obligations. Proof-based methods tend to be the most expressive and
general of all approaches to parameterized verification, capable of handling
protocols with complex data structures, arithmetic, and dynamic topology.
Their principal limitation is the burden they place on the user: finding and
expressing the right invariant requires deep insight into the protocol's
correctness argument, and scaling to large or complex protocols remains
demanding even with modern tool support.
